% GENERATED FILE. Edit canonical lessons, then run npm run generate:books.
% canonical-source-hash: fnv1a64:e3b71f337fd6d33c
% canonical-lessons: SA-C88-girl, SA-C88-youngman, SA-C88-friend, SA-C88-son, SA-C88-daughter

\chapter{Girls, Boys and Friends}
\label{ch:sa-girls-boys-friends}
\hlchaptermodality{88}

\begin{chapteropening}
\textbf{By the end of this chapter:} \emph{I can name a girl, a young man, a woman's friend, a son and a daughter.}

Complete the last lesson of chapter 88: i can name a girl, a young man, a woman's friend, a son and a daughter.
\end{chapteropening}

\section[bālikā]{\sk{बालिका} (bālikā) — a girl}

\label{lesson:SA-C88-girl}



\begin{quote}
Before the new one: say the Sanskrit for a hundred, then the Sanskrit for a thousand.
\end{quote}

\subsection*{You'll want to know: \sk{बालिका}}
\textbf{\sk{बालिका}} (\emph{bālikā}) — \textquotedblleft{}a girl\textquotedblright{}. \textbf{\sk{बालः}} is \textquotedblleft{}a boy\textquotedblright{}; \textbf{-\sk{इका}} makes the girl.

\begin{grammarlens}[title={What you've built}]
The boy word, made feminine.
\end{grammarlens}

\subsection*{Your turn}
\begin{itemize}
\raggedright
  \item \emph{Say it:} \emph{bālikā}
  \item \emph{Say it:} \emph{bālikā}, once more
  \item \emph{Say it:} say \emph{bālikā}
  \item \emph{Recall:} say the Sanskrit for a hundred, then the Sanskrit for a thousand, then say \emph{bālikā} again
\end{itemize}

\begin{tcolorbox}[breakable,colback=teal!4,colframe=teal!35!black,title={Before you move on}]
What does \sk{बालिका} mean? (\textquotedblleft{}A girl\textquotedblright{}.) Say it once more.
\end{tcolorbox}
\section[yuvakaḥ]{\sk{युवकः} (yuvakaḥ) — a young man}

\label{lesson:SA-C88-youngman}



\begin{quote}
Before the new one: say the Sanskrit for a thousand, then the Sanskrit for a girl.
\end{quote}

\subsection*{You'll want to know: \sk{युवकः}}
\textbf{\sk{युवकः}} (\emph{yuvakaḥ}) — \textquotedblleft{}a young man\textquotedblright{}, between a \textbf{\sk{बालः}} and a \textbf{\sk{वृद्धः}}.

\begin{grammarlens}[title={What you've built}]
Between boy and old man.
\end{grammarlens}

\subsection*{Your turn}
\begin{itemize}
\raggedright
  \item \emph{Say it:} \emph{yuvakaḥ}
  \item \emph{Say it:} \emph{yuvakaḥ}, once more
  \item \emph{Say it:} say \emph{yuvakaḥ}
  \item \emph{Recall:} say the Sanskrit for a thousand, then the Sanskrit for a girl, then say \emph{yuvakaḥ} again
\end{itemize}

\begin{tcolorbox}[breakable,colback=teal!4,colframe=teal!35!black,title={Before you move on}]
What does \sk{युवकः} mean? (\textquotedblleft{}A young man\textquotedblright{}.) Say it once more.
\end{tcolorbox}
\section[sakhī]{\sk{सखी} (sakhī) — a (female) friend}

\label{lesson:SA-C88-friend}



\begin{quote}
Before the new one: say the Sanskrit for a girl, then the Sanskrit for a young man.
\end{quote}

\subsection*{You'll want to know: \sk{सखी}}
\textbf{\sk{सखी}} (\emph{sakhī}) — \textquotedblleft{}a friend\textquotedblright{}, said of a woman or a girl; \textbf{\sk{मित्रम्}} is any friend.

\begin{grammarlens}[title={What you've built}]
A friend, in the feminine.
\end{grammarlens}

\subsection*{Your turn}
\begin{itemize}
\raggedright
  \item \emph{Say it:} \emph{sakhī}
  \item \emph{Say it:} \emph{sakhī}, once more
  \item \emph{Say it:} say \emph{mama sakhī}
  \item \emph{Recall:} say the Sanskrit for a girl, then the Sanskrit for a young man, then say \emph{sakhī} again
\end{itemize}

\begin{tcolorbox}[breakable,colback=teal!4,colframe=teal!35!black,title={Before you move on}]
What does \sk{सखी} mean? (\textquotedblleft{}A (female) friend\textquotedblright{}.) Say it once more.
\end{tcolorbox}
\section[putraḥ]{\sk{पुत्रः} (putraḥ) — a son}

\label{lesson:SA-C88-son}



\begin{quote}
Before the new one: say the Sanskrit for a young man, then the Sanskrit for a woman's friend.
\end{quote}

\subsection*{You'll want to know: \sk{पुत्रः}}
\textbf{\sk{पुत्रः}} (\emph{putraḥ}) — \textquotedblleft{}a son\textquotedblright{}, the everyday word beside \textbf{\sk{सूनुः}}.

\begin{grammarlens}[title={What you've built}]
A second word for son.
\end{grammarlens}

\subsection*{Your turn}
\begin{itemize}
\raggedright
  \item \emph{Say it:} \emph{putraḥ}
  \item \emph{Say it:} \emph{putraḥ}, once more
  \item \emph{Say it:} say \emph{mama putraḥ}
  \item \emph{Recall:} say the Sanskrit for a young man, then the Sanskrit for a woman's friend, then say \emph{putraḥ} again
\end{itemize}

\begin{tcolorbox}[breakable,colback=teal!4,colframe=teal!35!black,title={Before you move on}]
What does \sk{पुत्रः} mean? (\textquotedblleft{}A son\textquotedblright{}.) Say it once more.
\end{tcolorbox}
\section[putrī]{\sk{पुत्री} (putrī) — a daughter}

\label{lesson:SA-C88-daughter}



\begin{quote}
Before the new one: say the Sanskrit for a woman's friend, then the Sanskrit for a son.
\end{quote}

\subsection*{You'll want to know: \sk{पुत्री}}
\textbf{\sk{पुत्री}} (\emph{putrī}) — \textquotedblleft{}a daughter\textquotedblright{}, \textbf{\sk{पुत्रः}} ending in a long \textbf{\sk{ी}}: the everyday word beside \textbf{\sk{दुहिता}}.

\begin{grammarlens}[title={What you've built}]
Son, made feminine.
\end{grammarlens}

\subsection*{Your turn}
\begin{itemize}
\raggedright
  \item \emph{Say it:} \emph{putrī}
  \item \emph{Say it:} \emph{putrī}, once more
  \item \emph{Say it:} say \emph{mama putrī}
  \item \emph{Recall:} say the Sanskrit for a woman's friend, then the Sanskrit for a son, then say \emph{putrī} again
\end{itemize}

\begin{tcolorbox}[breakable,colback=teal!4,colframe=teal!35!black,title={Before you move on}]
What does \sk{पुत्री} mean? (\textquotedblleft{}A daughter\textquotedblright{}.) Say it once more.
\end{tcolorbox}
